% ============================================================
% Related work — spear phishing, susceptibility, graphs, methodology (target: C&S)
% ============================================================
\section{Related Work}
\label{sec:related}

The review follows the chain of the thesis: spear phishing as a violation of domain knowledge
(\ref{sec:rw-spear}), human and organizational susceptibility (\ref{sec:rw-suscept}), a graph signal in place of
content (\ref{sec:rw-graph}), and evaluation rigor (\ref{sec:rw-eval}), closing with the gaps (\ref{sec:rw-gap}).

\subsection{Spear phishing and BEC: why content is insufficient}
\label{sec:rw-spear}
Targeted attacks, such as spear phishing and whaling, alongside Business Email Compromise, differ from mass phishing in their personalization and context.
A recent survey in \emph{Computers \& Security} systematizes the techniques and gaps
of the area \citep{birthriya2025detection}. \citet{ho2019detecting} characterized lateral phishing from
compromised accounts across $113$~million emails and fixed an operational low-FPR threshold of $<\!4$~FP/million.
\citet{cidon2019high} targeted payload-free BEC in production, and BEC surveys trace its evolution and cost
\citep{atlam2023business}. What the effective methods share is a departure from content: per-sender behavioral
profiles \citep{gascon2018reading,stringhini2015thataint}, header and stylometry models
\citep{duman2016emailprofiler}, as well as analysis of relay infrastructure \citep{luo2024characterizing} all catch mail
from compromised legitimate accounts whose content reads as normal. In each case, maliciousness
is determined by the sender and the surrounding context rather than by what was written.

\subsection{The knowledge-side attack: OSINT and language models}
\label{sec:rw-attack}
Spear phishing draws its credibility from \emph{domain knowledge} built on open-source intelligence (OSINT).
Automated harvesting of such traces has been studied since early work \citep{edwards2017panning}. Yet \citet{dewan2014analyzing}
found that raw social and OSINT features added \emph{little} beyond content and URLs
\citep{dewan2014analyzing,dewan2017facebook}, a caution we take seriously: OSINT helps as a \emph{relation
constructor}, a shared list or affiliation, and treating it as a flat feature is a known pitfall. Language models then drastically lower the
cost of personalization. Automated spear-phishing campaigns match human experts in human-subject studies
\citep{bethany2024evaluating}, and realistic OSINT traces can be generated synthetically without subjects
\citep{brown2025web}. Once content can be
generated perfectly, the defense must exploit relations rather than style.

\subsection{Human and organizational susceptibility}
\label{sec:rw-suscept}
The human factor explains why this hard-to-detect tail is so effective. Susceptibility to spear phishing depends on
personality traits and context \citep{halevi2015spearphishing}, and time pressure together with persuasion cues
(authority, urgency) reduces detectability \citep{butavicius2022why}. Awareness and decision style
moderate detection ability \citep{sturman2025security}. The implication for technical defense is significant:
because susceptibility is a function of \emph{role and relation, alongside context}, (a)~a content filter is
insufficient, and (b)~it can be \emph{measured} structurally, which we do via an organizational
susceptibility metric (Section~\ref{sec:susceptibility}).

\subsection{The graph signal: from GNN foundations to knowledge graphs}
\label{sec:rw-graph}
The two invariants rest on a mature apparatus of graph learning, organized here into six threads that indicate
what we borrow and what is missing for a \emph{domain knowledge graph} of email.

On the foundations of inductive learning, modern graph networks learn node representations through
neighborhood aggregation: the \emph{inductive} GraphSAGE \citep{hamilton2017inductive} (generalization to nodes
unseen during training) and the attention-based GAT \citep{velickovic2018graph} establish the paradigm; our \emph{structural} (rather than identity-based) node features are a deliberately
lightweight variant of it, to preserve inductive generalization \citep{xu2020inductive,hao2020inductive}. Link-prediction and scalable-embedding tasks
\citep{zhang2018link,zhang2018scalable} provide heads for classifying $(s\!\to\!v)$ events.

The formalism for fusing multiple relations over \emph{one}
set of accounts is the \emph{multiplex} network
\citep{kivela2014multilayer,dedomenico2013mathematical,magnani2013combinatorial,jeude2019multilayer}.
Heterogeneous GNNs learn to fuse relation types: HAN \citep{wang2019heterogeneous}, HGT \citep{hu2020heterogeneous},
MAGNN \citep{fu2020magnn}, and attributed-network embedding \citep{christophides2021overview,cen2019representation}.
This provides the machinery for joining layers, yet \emph{cross-layer (in)consistency as a security signal}
has not been operationalized for corporate email.

The second invariant requires a model of time: the memory-based TGN
\citep{rossi2020temporal}, the attention-based DySAT \citep{sankar2020dysat}, EvolveGCN \citep{pareja2020evolvegcn},
ROLAND \citep{you2022roland}, and streaming and spatio-temporal GNNs
\citep{yu2023streaming,wang2022tsgn,nayebi2025spatiotemporal,behrouz2022anomaly} capture graph evolution. Our
TGN-lite is a minimal realization of this paradigm by design, allowing the time-shuffle control to isolate the
\emph{mechanism} rather than the capacity of the architecture.

The closest apparatus comes from fraud detection robust to
\emph{camouflage} and \emph{class imbalance}: CARE-GNN \citep{dou2020enhancing}, GraphConsis \citep{liu2020alleviating},
PC-GNN \citep{liu2021pick}, H2-FDetector \citep{shi2022h2fdetector}, GraphSMOTE \citep{zhao2021graphsmote},
distribution-shift correction \citep{gao2023alleviating,tang2022rethinking}, and the detection of money laundering
and transactional fraud \citep{weber2019antimoney,li2022ttagn,rayana2015collective}. The broader graph-anomaly
literature \citep{ding2019deep,park2020unsupervised,liu2021anomaly,ma2021comprehensive,zhang2021mcgc}, (including the multiplex setting \citep{li2025umgad}), as well as predicting user dynamics and roles
\citep{kumar2019predicting,wu2020who,li2022pdgnn,mcauley2015image} provides techniques, yet operates on graphs of
\emph{transactions/reviews} and has yet to address \emph{communication with domain knowledge}.

Since an attacker can \emph{fabricate} edges
(our fabrication sweep), the literature on graph structure attacks
\citep{zugner2018adversarial,zugner2019adversarial,dai2022towards} and defenses
\citep{zhang2020gnnguard,jin2020graph,cai2021structural,wang2023label} is relevant. Contrastive and
self-supervised learning \citep{hassani2020contrastive,wang2021selfsupervised,singh2024node,yan2024topological}
offers a label-free signal, relevant given the absence of real BEC labels.

Graph-based phishing detection has so far
targeted mainly \emph{web pages} and \emph{transactions}: graph-theoretic and visual approaches for web pages
\citep{tan2020graph,lin2021phishpedia,abdelnabi2020visualphishnet}, attention over the email graph
\citep{safran2025phishinggnn}, edge-feature GraphSAGE \citep{lo2021egraphsage}, propagation over URL--IP graphs
\citep{guo2025efficient,li2022pdgnn}, chain/blockchain-based detection
\citep{daluwatta2022cgraph,edirimannage2022phishchain}, and the detection of phishing and accounts in social
media \citep{aggarwal2013phishari,shu2017user,koide2024chatspamdetector,song2025insight,dube2024building}. In
parallel, cyber threat intelligence (CTI) knowledge graphs structure intelligence
\citep{cheng2024ctinexus,fieblinger2024actionable,li2021siege}. \emph{Cross-view (in)consistency} as an anomaly
signal is known \citep{liu2020alleviating}, but to our knowledge it has not been operationalized as a \emph{domain
knowledge graph} of corporate email for security.

\subsection{Evaluation rigor: pitfalls and leak control}
\label{sec:rw-eval}
A central caution is that ML detectors in security are easily inflated by leaks, unfair splitting, and headline
AUC in place of operational metrics \citep{arp2022dos}; time-aware evaluation (TESSERACT) shows how temporal
leakage distorts results \citep{pendlebury2019tesseract}. These works motivate our leak-aware protocol and
predictivity probe, which ask not how high the AUC is but whether the result transfers beyond a single corpus and
whether it is an artifact of the generator. We use the same standard to define a method's effectiveness.

\subsection{Gaps and positioning of the work}
\label{sec:rw-gap}
Several gaps emerge from the review. Spear phishing detection remains dominated by the content signal, while
the targeted tail bypasses it by definition, and there is no unified relational treatment for lateral/BEC or external
impersonation. Graph approaches target the web and transactions rather than the domain knowledge graph of
communication, and organizational susceptibility is not measured structurally. Evaluations also rarely control for leaks or test transfer across topologies, so effectiveness tends to be
ill-defined. The present work
addresses these gaps: a domain knowledge graph, its two detection invariants, a susceptibility metric, and a
leak-aware methodology that defines effectiveness.
